\chapter{LaTeX Examples}
\label{cha:latex_examples}

Este capítulo deve conter uma síntese inicial de trabalho relacionado (como resumo de estado da arte) e com foco nos objectivos e contribuições da dissertação. O capítulo deve demonstrar que o aluno procedeu e uma recolha bibliográfica para análise, tratamento e classificação preliminar, devendo esta ser suficientemente representativa das leituras preparatórias para o prosseguimento da elaboração da dissertação. Esta síntese deve ser apresentada com as respectivas referências bibliográficas (correctamente apresentadas na Bibliografia).

Usar no capítulo fonte palatino 12 com um espaçamento e meio, esperando-se uma dimensão indicativa de cerca de 20 páginas.


Para fazer citações, deverá usar-se a chave da referência no ficheiro BibTeX. Se for uma única referência~\cite{Artho04}, usar um ``\verb!~!'' para ligar o \verb!\cite{...}! à palavra que o precede (\ldots\verb!referência~\cite{Artho04}!).  Caso queira fazer múltiplas citações~\cite{Shavit95,Silberschatz06,Moss85}, deverá agrupá-las dentro de um úinico \verb!\cite{...}!.

Note que o ficheiro de bibliografia pode ter tantas entradas quantas quiser. Apenas aquelas cuja chave seja referenciada no texto é que serão incluidas na listagem de bibliografia.


Footnotes\footnote{This is a simple footnote.} will be numbered and shown in the bottom of the page.


A Tabela~\ref{tab:hla:results} ilustra alguns conceitos importantes associados à contrução de tabelas:
\begin{asparaenum}[i)]
	\item Não usar linhas verticais;
	\item A legenda deve ficar por cima da tabela;
	\item Usar as macros \verb!\toprule!, \verb!\midrule! e \verb!\bottomrule! para fazer a linha horizontal superior, interiores e inferior, respectivamente.
\end{asparaenum}
 
\begin{table}[ht]
	\caption{Test results summary.}
	\label{tab:hla:results}
\centering
\begin{tabular}{lccccc}
	\toprule
	\multicolumn{1}{c}{\textbf{Test}} 	& \textbf{Anomalies}	& \textbf{Warnings}	& \textbf{Correct} 	& \textbf{Categories}		& \textbf{Missed} \\
	\midrule
\cite{Beckman08}~Connection 	& 2 & 2	& 1	& \emph{C}				& 1 \\
\cite{Artho03}~Coordinates'03 	& 1	& 4	& 1	& \emph{2B, 1C}			& 0 \\
\cite{Artho03}~Local Variable	& 1	& 2	& 1	& \emph{A}				& 0 \\
\cite{Artho03}~NASA				& 1	& 1	& 1	& ---					& 0 \\
\cite{Artho04}~Coordinates'04	& 1	& 4	& 1	& \emph{3C}				& 0 \\
\cite{Artho04}~Buffer			& 0	& 7	& 0	& \emph{2A, 1B, 2C, 2D}	& 0 \\
\cite{Artho04}~Double-Check		& 0	& 2	& 0	& \emph{1A, 1B}			& 0 \\
\cite{Flanagan04}~StringBuffer	& 1	& 0	& 0	& ---					& 1 \\
\cite{Praun03}~Account			& 1	& 1	& 1	& ---					& 0 \\
\cite{Praun03}~Jigsaw			& 1	& 2	& 1	& \emph{C}				& 0 \\
\cite{Praun03}~Over-reporting	& 0	& 2	& 0	& \emph{1A, 1C}			& 0 \\
\cite{Praun03}~Under-reporting	& 1	& 1	& 1	& ---					& 0 \\
\cite{IBM-Rep}~Allocate Vector	& 1	& 2	& 1	& \emph{C}				& 0 \\
Knight Moves					& 1	& 3	& 1	& \emph{2B}				& 0 \\
	\midrule
	\textbf{Total}			& \textbf{12}		& \textbf{33}		& \textbf{10}			& \textbf{5A, 6B, 10C, 2D}	& \textbf{2} \\
	\bottomrule
\end{tabular}
\end{table}


As figuras a inserir no documento deverão ser de qualidade, preferencialmente em formato vectorial (PDF vectorial) e não em \emph{bitmap} (PNG, JPG, etc). As imagens \emph{bitmap} (Figura~\ref{fig:Figuras_Tree_silhouettes-bitmap}) não escalam bem e têm reflexos negativos na qualidade do seu docuemnto.  Pelo contrário, as imagens \emph{vectoriais} {Figura~\ref{fig:Figuras_Tree_silhouettes-vectorial}} escalam muito tanto quanto o necessário sem degradar a qualidade da imagem.

Só deve usar \emph{screenshots} se não tive mesmo nenhuma alternativa.  Em vez e gerar um \emph{screenshot}, tente usar uma impressora virtual PDF e imprimir para um ficheiro PDF. Regra geral obterá um PDF vetorial. Mesmo que o seu PDF contenha imagens, elas terão sempre qualidade maior ou igual à que obteria com um \emph{screenshot}.


\begin{figure}[htbp]
	\centering
	\includegraphics[height=1in]{snowman-bitmap}
	\includegraphics[height=3in]{snowman-bitmap}
	\includegraphics[height=6in]{snowman-bitmap}
	\caption{Imagem em formato \emph{bitmap} (JPG)}
	\label{fig:Figuras_Tree_silhouettes-bitmap}
\end{figure}

\begin{figure}[htbp]
	\centering
	\includegraphics[height=1in]{snowman-vectorial}
	\includegraphics[height=3in]{snowman-vectorial}
	\includegraphics[height=6in]{snowman-vectorial}
	\caption{Imagem em formato PDF vectorial}
	\label{fig:Figuras_Tree_silhouettes-vectorial}
\end{figure}

Pode usar o pacote \emph{subfigure} para agragar várias figuras numa única. Poderá assim referenciar o conjunto~\ref{fig:figura-completa}, a priemira delas~\ref{fig:novelo} ou a segunda~\ref{fig:nuvem}.


\begin{figure}[htbp]
	\centering
    \subfigure[Novelo de lã] {
		\label{fig:novelo}
		\includegraphics[height=1in]{knitting-vectorial}
    }
\qquad\qquad
    \subfigure[Tempestade com neve] {
		\label{fig:nuvem}
		\includegraphics[height=1in]{snowstorm-vectorial}
    }
  \caption{Exemplo de utilização de \emph{subfigure}}
  \label{fig:figura-completa}
\end{figure}


Para incluir listagens de código no seu documento, deverá incluir o pacote \emph{listings} e depois usar o ambiente \emph{lstlisting}, como exemplificado na Listagem~\ref{lst:HelloWorld}.

\lstset{language=Java, caption=Hello World, label=lst:HelloWorld}
\begin{lstlisting}
/** 
 * The HelloWorldApp class implements an application that
 * simply prints "Hello World!" to standard output.
 */
class HelloWorldApp {
    public static void main(String[] args) {
        System.out.println("Hello World!"); // Display the string.
    }
}
\end{lstlisting}

\section{Equações}

O LaTeX é uma ferramenta poderosa para escrever em estilo matemático. Permite inserir fórmulas no meio do texto como por exemplo esta: $ax^2 + bx + c = 0$. Também permite que as fórmulas sejam destacadas numa linha separada e centradas na página 
$$x = \frac{-b \pm \sqrt{b^2-4ac}}{2a}$$
\[x = \frac{-b \pm \sqrt{b^2-4ac}}{2a}\]
ou numeradas 
\begin{equation}
aaa
\label{eq:1}
\end{equation}
que depois pode ser referida no texto como sendo a equação~\ref{eq:1}
$$\begin{array}{l}
aa
\end{array}
$$

\begin{eqnarray}
a\\
b\\
c\\
\end{eqnarray}

Contributors for the examples:\\
\begin{tabular}{l}
João Lourenço (Joao.Lourenco@di.fct.unl.pt)\\
João Seco (Joao.Seco@di.fct.unl.pt)\\
Luís Russo (Luis.Russo@di.fct.unl.pt)\\
Vitor Duarte (Vitor.Duarte@di.fct.unl.pt)\\
Ricardo Dias (rjfd@di.fct.unl.pt)\\
\end{tabular}